% ==== XII ====
\yoschapter{XII}{Other Representation Theorems in Linear Spaces}
\label{chap:XII}

In this chapter, we shall prove three representation theorems in linear spaces. The first one, the Krein-Milman theorem says that a non-void convex compact subset $K$ of a locally convex linear topological space is equal to the closure of the convex hull of the extremal points of $K$. The other two theorems concern the representations of a vector lattice as point functions and as set functions. 


\yossection{1}{Extremal Points. The Krein-Milman Theorem}
\label{sec:XII.1}

\begin{definitionu} Let $K$ be a subset of a real- or complex-linear space $X$. A non-void subset $M \subseteq K$ is said to be an extremal subset of $K$, if a proper convex combination $\alpha k_{1} + (1 - \alpha) k_{2}, 0 < \alpha < 1$, of two points $k_{1}$ and $k_{2}$ of $K$ lies in $M$ only if both $k_{1}$ and $k_{2}$ are in $M$. An extremal set of $K$ consisting of just one point is called an extremal point of $K$. 
\end{definitionu}

\begin{examplen} In a three dimensional Euclidean space, the surface of a solid sphere is an extremal subset of the sphere, and every point of the surface is an extremal point of the surface.

\end{examplen}
\begin{theoremn}[\textsc{Krein}-\textsc{Milman}] A non-void compact convex subset $K$ of a locally convex linear topological space $X$ has at least one extremal point.

\end{theoremn}
\begin{proof}
 The set $K$ is itself an extremal set of $K$. Let $M$ be the totality of compact extremal subsets $M$ of $K$. Order $M$ by inclusion relation. It is easy to see that, if $M_{1}$ is a linearly ordered subfamily of $M$, the non-void set $\bigcap\limits_{M\in M_{1}}M$ is a compact extremal subset of $K$ which is a lower bound for the subfamily $M_{1}$. 

Thus, by Zorn's lemma, $\mathfrak{M}$ contains a minimal element $M_0$. Suppose that $M_0$ contains two distinct points $x_0$ and $y_0$. Then there exists a continuous linear functional $f$ on $X$ such that $f(x_0) \neq f(y_0)$. We may assume that $\operatorname{Re} f(x_0) \neq \operatorname{Re} f(y_0)$. $M_0$ being compact, the set $M_1 = \{x \in M_0; \operatorname{Re} f(x) = \inf\limits_{y \in M_0} \operatorname{Re} f(y)\}$ is a proper subset of $M_0$. On the other hand, if $k_1$ and $k_2$ are points of $K$ such that $\alpha k_1 + (1 - \alpha) k_2 \in M_1$ for some $\alpha$ with $0 < \alpha < 1$, then, by the extremal property of $M_{0}$, $k_{1}$ and $k_{2}$ both $\in M_{0}$. It follows from the definition of $M_{1}$ that $k_{1}$ and $k_{2}$ both $\in M_{1}$. Hence $M_{1}$ is a closed extremal subset properly contained in $M_{0}$. Since $M_{0}$ is a minimal element of $M$, we have arrived at a contradiction. Therefore $M_{0}$ must consist of only one point which is thus an extremal point of $K$. 
\end{proof}
\begin{corollaryu} Let $K$ be a non-void compact convex subset of a locally convex real linear topological space $X$. Let $E$ be the totality of the extremal points of $K$. Then $K$ coincides with the smallest closed set containing every convex combination $\sum\limits_{i} \alpha_{i} e_{i} (\alpha_{i} \geqq 0, \sum\limits_{i} \alpha_{i} = 1)$ of points $e_{i} \in E$, i.e., $K$ is equal to the closure of the convex hull Conv $(E)$ of $E$. 
\end{corollaryu}

\begin{proof}
 The inclusion $E \subseteq K$ and the convexity of $K$ imply that $\Conv(E)^a$ is contained in the compact set $K$. Suppose that there exists a point $k_0$ contained in $(K - \Conv(E)^a)$. We can then take a point $c \in \Conv(E)^a$ so that $(k_0 - c) \bar{\in} (\Conv(E)^a - c)$. The set $(\Conv(E)^a - c)$ being compact convex and $\ni 0$, there exists, by Theorem~\ref{thm:IV.6.Theorem.3prime} in Chapter~\ref{chap:IV}, \ref{sec:IV.6}, a continuous real linear functional $f$ on $X$ such that 
\[
f (k _{0} - c) > 1 \text{ and } f (k - c) \leqq 1 \text{ for } (k - c) \in (\Conv (E) ^{a} - c). \]

Let $K_{1} = \{x \in K; f(x) = \sup\limits_{y \in K} f(y)\}$. Then, since $k_{0} \in K$, the set $K_{1} \cap E$ must be void. Moreover, since $K$ is compact, $K_{1}$ is a closed extremal subset of $K$. On the other hand, any extremal subset of $K_{1}$ is also an extremal subset of $K$ and hence any extremal point of $K_{1}$, which surely exists by the preceding Theorem, is also an extremal point of $K$. Since $K_{1} \cap E$ is void, we have arrived at a contradiction. 
\end{proof}
\begin{remarku} The above Theorem and the Corollary were first proved by M. \textsc{Krein}-D. \textsc{Milman} \cite{krein165}. The proof given above is adapted from J. L. \textsc{Kelley} \cite{kelley147}. It is to be noted that for the unit sphere $S = \{x \in X; \| x \| \leqq 1\}$ in a Hilbert space $X$ the extremal points of $S$ are precisely those on the surface of $S$, i.e., those of norm $1$. This we easily see from \eqref{eq:I.5.1} in Chapter~\ref{chap:I}, \ref{sec:I.5}. For applications of the notion of the extremal points to concrete function spaces, see, e.g., K. \textsc{Hoffman} \cite{hirsch108}. 
\end{remarku}

A simple example. Let $C[0,1]$ be the space of real-valued continuous functions $x(t)$ defined on $[0,1]$ normed by $x = \max\limits_t |x(t)|$. The dual space $X = C[0,1]'$ is the space of real Baire measures on $[0,1]$ of bounded total variations. The unit sphere $K$ of $X$ is compact in the weak* topology of $X$ (see Theorem~\hyperref[thm:AV.1.Theorem.1]{$1$} in the Appendix to Chapter~\ref{chap:IV}). It is easy to see that extremal points of $K$ are in one-to-one correspondence with the linear functionals $f_{t_0} \in X$ of the form $\langle x, f_{t_0} \rangle = x(t_0), t_0 \in [0,1]$. The above Corollary says that any linear functional $f \in X$ is given as the weak* limit of the functionals of the form 
\[
\sum\limits_{j=1} ^{n} \alpha_{j} x (t _{j}), \quad \text{ where } \quad \alpha_{j} > 0, \sum\limits_{j=1} ^{n} \alpha_{j} = 1 \quad \text{ and } \quad t _{j} \in [ 0, 1 ].
\]

Recently, G. \textsc{Choquet} \cite{choquet33} proved a more precise result: If $X$ is a metric space, then $E$ is a $G_{\delta}$ set and, for every $x \in K$, there exists a non-negative Baire measure $\mu_x(B)$ defined for Baire sets $B$ of $X$ such that $\mu_x(X - E) = 0$, $\mu_x(E) = 1$ and $x = \int_E y \mu_x(\dd y)$. As for the uniqueness of $\mu_x$ and further literature, see G. \textsc{Choquet} and P. A. \textsc{Meyer} \cite{choquet34}.\yosstarfootnote{See also Supplementary Notes, p.~\pageref{notes:XII}.} 


\yossection{2}{Vector Lattices}
\label{sec:XII.2}

The notion of ``positivity'' in concrete function spaces is very important, in theory as well as in application. A systematic abstract treatment of the ``positivity'' in linear spaces was introduced by F. \textsc{Riesz} \cite{rickart277}, and further developed by H. \textsc{Freudenthal} \cite{freudenthal67}, G. \textsc{Birkhoff} \cite{birkhoff17} and many other authors. These results are called the theory of vector lattice. We shall begin with the definition of the vector lattice. 

\begin{definition}\label{def:XII.2.Definition.1} $A$ real linear space $X$ is said to be a vector lattice if $X$ is a lattice by a partial order relation $x \leqq y$ satisfying the conditions: 
\begin{equation}\label{eq:XII.2.1}
x \leqq y \text{ implies } x + z \leqq y + z,
\end{equation}
\begin{equation}\label{eq:XII.2.2}
x \leqq y \text{ implies } \alpha x \leqq \alpha y \text{ (or } \alpha x \geqq \alpha y\text{) for every } \alpha \geqq 0 \text{ (or } \alpha \leqq 0\text{)}.
\end{equation}
\end{definition}


\begin{proposition}\label{prop:XII.2.Proposition.1} If, in a vector lattice $X$, we define 
\begin{equation}\label{eq:XII.2.3}
x ^{+} = x \vee 0 \text{ and } x ^{-} = x \wedge 0,
\end{equation}
then we have 
\begin{equation}\label{eq:XII.2.4}
x \vee y = (x - y) ^{+} + y, x \wedge y = - ((- x) \vee (- y)).
\end{equation}
\end{proposition}

\begin{proof}
 The one-one mappings $x \rightarrow x + z$ and $x \rightarrow \alpha x (\alpha > 0)$ of $X$ onto $X$ both preserve the partial order in $X$. 
\end{proof}

\begin{examplen} The totality $A(S, \mathfrak{B})$ of real-valued, $\sigma$-additive set functions $x(B)$ defined and finite on a $\sigma$-additive family $(S, \mathfrak{B})$ of sets $B \subseteq S$ is a vector lattice by
\[
(x + y) (B) = x (B) + y (B), (\alpha x) (B) = \alpha x (B)
\]
and the partial order $x \leqq y$ defined by $x(B) \leqq y(B)$ on $\mathfrak{B}$. In fact, we have, in this case, 
\begin{equation}\label{eq:XII.2.5}
x ^{+} (B) = \sup\limits _{N \subseteq B} x (N) = \text{ the positive variation } \overline {{V}} (x; B) \text{ of } x \text{ on } B.
\end{equation}

\end{examplen}
\begin{proof}
 We have to show that $\overline{V}(x; B) = (x \vee 0)(B)$. It is clear that $\overline{V}(x; B) \geqq 0$ and $x(B) \leqq \overline{V}(x; B)$ on $\mathfrak{B}$. If $0 \leqq y(B)$ and $x(B) \leqq y(B)$ on $\mathfrak{B}$, then, for any $N \subseteq B$, $y(B) = y(N) + y(B - N) \geqq x(N)$ and so $y(B) \geqq \overline{V}(x; B)$ on $\mathfrak{B}$. 
\end{proof}

\begin{proposition}\label{prop:XII.2.Proposition.2} In a vector lattice $X$, we have 
\begin{align}
x \underset {\wedge} {\vee} y + z& = (x + z) \underset {\wedge} {\vee} (y + z), \label{eq:XII.2.6} \\
\alpha (x \underset {\wedge} {\vee} y)& = (\alpha x) \underset {\wedge} {\vee} (\alpha y) \quad \text{ for } \quad \alpha > 0, \label{eq:XII.2.7} \\
\alpha (x \underset {\wedge} {\vee} y)& = (\alpha x) \underset {\vee} {\wedge} (\alpha y) \quad \text{ for } \quad \alpha < 0, \label{eq:XII.2.8} \\
x \wedge y& = - (- x) \vee (- y), x ^{-} = - (- x) ^{+}, x ^{+} = - (- x) ^{-}. \label{eq:XII.2.9}
\end{align}
\end{proposition}

\begin{proof}
 Clear from \eqref{eq:XII.2.1} and \eqref{eq:XII.2.2}. 
\end{proof}

\begin{corollaryu} 
\begin{equation}\label{eq:XII.2.10}
x+y=x\vee y+x\wedge y,\qquad\text{in particular, }x=x^++x^-.
\end{equation}
\end{corollaryu}

\begin{proof}
\[
x\vee y-x-y=0\vee(y-x)-y=(-y)\vee(-x)=-(y\wedge x).
\]
\end{proof}



\begin{proposition}\label{prop:XII.2.Proposition.3} We have 
\begin{align}
x \underset {\wedge} {\vee} y& = y \underset {\wedge} {\vee} x (\text{ commutativity }), \label{eq:XII.2.11} \\
x \vee (y \vee z)& = (x \vee y) \vee z = \sup (x, y, z) \Big \} _ {\text{ (associativity) }} \label{eq:XII.2.12} \\
x \wedge (y \wedge z)& = (x \wedge y) \wedge z = \inf (x, y, z) \left\} ^ {(\text{ associativity })}, \right. \label{eq:XII.2.13} \\
&\quad (x \wedge y) \vee z = (x \vee z) \wedge (y \vee z) \quad (\text{ distributivity }), \label{eq:XII.2.14} \\
&\quad (x \vee y) \wedge z = (x \wedge z) \vee (y \wedge z) \quad (\text{ distributivity }). \label{eq:XII.2.15}
\end{align}
\end{proposition}


\begin{proof}
 We have only to prove the distributivity. It is clear that $(x \wedge y) \vee z \leqq x \vee z, y \vee z$. Let $w \leqq x \vee z, y \vee z$. Then $w \leqq x \vee z = x + z - x \wedge z$ and so $x + z \geqq w + x \wedge z$. Similarly we have $y + z \geqq w + y \wedge z$. Hence 
\[
\begin{aligned}
&w + (x \wedge z) \wedge (y \wedge z) = (w + x \wedge z) \wedge (w + y \wedge z) \leqq (x + z) \wedge (y + z) = x \wedge y + z,
\end{aligned}
\]
and so 
\[
w \leqq (x \land y) + z - (x \land y) \land z = (x \land y) \lor z. \]

We have thus proved \eqref{eq:XII.2.14}. \eqref{eq:XII.2.15} is proved by substituting $-x$, $-y$, $-z$ for $x$, $y$, $z$, respectively in \eqref{eq:XII.2.14}. 
\end{proof}

\begin{remarkn}[Remark $1$] In a lattice, the distributive identity 
\begin{equation}\label{eq:XII.2.16}
x \wedge (y \vee z) = (x \wedge y) \vee (x \wedge z)
\end{equation}
\end{remarkn}

does not hold in general. The modular identity: 
\begin{equation}\label{eq:XII.2.17}
x \leqq z \text{ implies } x \vee (y \wedge z) = (x \vee y) \wedge z
\end{equation}
is weaker than the distributive identity \eqref{eq:XII.2.15}. Let $G$ be a group. Then the totality of invariant subgroups $N$ of $G$ constitutes a modular lattice, that is, a lattice satisfying the modular identity \eqref{eq:XII.2.17}, if we define $N_{1} \vee N_{2}$ and $N_{1} \wedge N_{2}$ as the invariant subgroup generated by $N_{1}$ and $N_{2}$ and the invariant subgroup $N_{1} \cap N_{2}$, respectively. 

\begin{remarkn}[Remark $2$] $A$ typical example of a distributive lattice is given by the Boolean algebra; a distributive lattice $B$ is called a Boolean algebra if it satisfies the conditions: \begin{enumerate}[(i)]
\item there exist elements $I$ and $0$ such that $0 \leqq x \leqq I$ for every $x \in B$,
\item for any $x \in B$, there exists a uniquely determined complement $x' \in B$ which satisfies $x \vee x' = I$, $x \wedge x' = 0$.
\end{enumerate}
 The totality $B$ of subsets of a fixed set is a Boolean algebra by defining the partial order in $B$ by the inclusion relation. 
\end{remarkn}

\begin{proposition}\label{prop:XII.2.Proposition.4} We define, in a vector lattice $X$, the absolute value 
\begin{equation}\label{eq:XII.2.18}
| x | = x \vee (- x).
\end{equation}
\end{proposition}

Then 
\begin{align}
| x |& \geqq 0, \text{ and } | x | = 0 \text{ iff } x = 0, \label{eq:XII.2.19} \\
| x + y |& \leqq | x | + | y |, | \alpha x | = | \alpha | | x |. \label{eq:XII.2.20}
\end{align}

\begin{proof}
 We have 
\begin{equation}\label{eq:XII.2.21}
x ^{+} \wedge (- x ^{-}) = x ^{+} \wedge (- x) ^{+} = 0. \end{equation}

For, by $0 = x - x = x \vee (-x) + x \wedge (-x) \geqq 2(x \wedge (-x))$ and the distributive identity \eqref{eq:XII.2.14}, $0 = (x \wedge (-x)) \vee 0 = x^{+} \wedge (-x)^{+}$. Thus 
\[
x ^{+} - x ^{-} = x ^{+} + (- x) ^{+} = x ^{+} \vee (- x) ^{+}.
\]

On the other hand, from $x \vee (-x) \geqq x \wedge (-x) \geqq -((-x) \vee x)$, we have $x \vee (-x) \geqq 0$ and so by \eqref{eq:XII.2.21} 
\[
x \vee (- x) = (x \vee (- x)) \vee 0 = x ^{+} \vee (- x) ^{+} = x ^{+} + (- x) ^{+}.
\]

We have thus proved 
\begin{equation}\label{eq:XII.2.22}
x ^{+} - x ^{-} = x ^{+} + (- x) ^{+} = x ^{+} \vee (- x) ^{+} = x \vee (- x).
\end{equation}

We now prove \eqref{eq:XII.2.19} and \eqref{eq:XII.2.20}. If $x = x^{+} + x^{-} \neq 0$, then $x^{+}$ or $(-x^{-})$ is $> 0$ so that $|x| = x^{+} \vee (-x^{-}) > 0$. $|\alpha x| = (\alpha x) \vee (-\alpha x) = |\alpha | (x \vee (-x)) = |\alpha | |x|$. From $|x| + |y| \geqq x + y, -x - y$, we obtain $|x| + |y| \geqq (x + y) \vee (-x - y) = |x + y|$. 
\end{proof}

\begin{remarku} The decomposition $x = x^{+} + x^{-}$, $x^{+} \wedge (-x^{-}) = 0$, is called the Jordan decomposition of $x; x^{+}, x^{-}$ and $|x|$ correspond to the positive variation, the negative variation and the total variation, of a function $x(t)$ of bounded variation, respectively. 
\end{remarku}

\begin{proposition}\label{prop:XII.2.Proposition.5} For any $y \in X$, we have 
\begin{equation}\label{eq:XII.2.23}
\left| x - x _{1} \right| = \left| x \vee y - x _{1} \vee y \right| + \left| x \wedge y - x _{1} \wedge y \right|.
\end{equation}
\end{proposition}

\begin{proof}
 We have 
\[ \begin{aligned}
&| a - b |  = (a - b) ^{+} - (a - b) ^{-} = a \vee b - b - (a \wedge b - b) = a \vee b - a \wedge b.
\end{aligned}
\]

Hence the right side of \eqref{eq:XII.2.23} is, by \eqref{eq:XII.2.10}, \eqref{eq:XII.2.14} and \eqref{eq:XII.2.15}, 
\[
\begin{aligned}
&= (x \lor y) \lor x _{1} - (x \lor y) \land (x _{1} \lor y) + (x \land y) \lor (x _{1} \land y) - (x \land y) \land x _{1}
\\
&= (x \lor x _{1}) \lor y - (x \land x _{1}) \lor y + (x \lor x _{1}) \land y - (x \land x _{1}) \land y
\\
&= x \lor x _{1} + y - (x \land x _{1} + y) = x \lor x _{1} - x \land x _{1} = | x - x _{1} |.
\end{aligned}
\]
\end{proof}

\begin{definition}\label{def:XII.2.Definition.2} $A$ sequence $\{x_{n}\}$ of a vector lattice $X$ is said to $O$-converge to an element $x \in X$, in symbol, $\Olim_{n \to \infty} x_{n} = x$ if there exists a sequence $\{w_{n}\}$ such that $|x - x_{n}| \leqq w_{n}$ and $w_{n} \downarrow 0$. Here $w_{n} \downarrow 0$ means that $w_{1} \geqq w_{2} \geqq \cdots$ and $\bigwedge\limits_{n \geqq 1} w_{n} = 0$. The \Olim $x_{n}$, if it exists, is uniquely determined. For, let \Olim $x_{n} = x$ and \Olim $x_{n} = y$. Then $|x - x_{n}| \leqq w_{n}$, $w_{n} \downarrow 0$ and $|y - x_{n}| \leqq u_{n}, u_{n} \downarrow 0$. Thus $|x - y| \leqq |x - x_{n}| + |x_{n} - y| \leqq w_{n} + u_{n}$ and $(w_{n} + u_{n}) \downarrow 0$ as may be seen from $w_{n} + \bigwedge\limits_{n \geqq 1} u_{n} \geqq \bigwedge\limits_{n \geqq 1} (w_{n} + u_{n})$. This proves that $x = y$. 
\end{definition}

\begin{proposition}\label{prop:XII.2.Proposition.6} With respect to the notion of \Olim, the operations $x + y, x \vee y$ and $x \wedge y$ are continuous in $x$ and $y$. 
\end{proposition}

\begin{proof}
 Let $\Olim_{n \to \infty} x_n = x$, $\Olim_{n \to \infty} y_n = y$. Then $|x + y - x_n - y_n| \leqq |x - x_n| + |y - y_n|$ implies $\Olim_{n \to \infty} (x_n + y_n) = x + y$. By \eqref{eq:XII.2.23}, we have 
\[
\begin{aligned}
&| x \underset {\wedge} {\vee} y - x _{n} \underset {\wedge} {\vee} y _{n} |  \leqq | x \underset {\wedge} {\vee} y - x _{n} \underset {\wedge} {\vee} y | + | x _{n} \underset {\wedge} {\vee} y - x _{n} \underset {\wedge} {\vee} y _{n} | \leqq | x - x _{n} | + | y - y _{n} |,
\end{aligned}
\]
and so \Olim $(x_{n} \wedge^{V} y_{n}) = (\Olim_{n \to \infty} x_{n}) \wedge^{V} (\Olim_{n \to \infty} y_{n})$. 
\end{proof}

\begin{remarku} We have 
\[
\Olim _{n \to \infty} \alpha \cdot x _{n} = \alpha \cdot \Olim _{n \to \infty} x _{n}.
\]
\end{remarku}

But, in general, 
\[
\Olim _{n \to \infty} \alpha_{n} x \neq \left(\lim\limits _{n \to \infty} \alpha_{n}\right) x.
\]

The former relation is clear from $\left|\alpha x - \alpha x_{n}\right| = \left|\alpha\right| \left|x - x_{n}\right|$. The latter inequality is proved by the following counter example: We introduce in the two-dimensional vector space the lexicographic partial order in which $(\xi_{1}, \eta_{1}) \geqq (\xi_{2}, \eta_{2})$ means that either $\xi_{1} > \xi_{2}$ or $\xi_{1} = \xi_{2}$, $\eta_{1} \geqq \eta_{2}$. It is easy to see that we obtain a vector lattice. We have, in this lattice, 
\[
n ^{-1} (1, 0) \geqq (0, 1) > 0 = (0, 0) (n = 1, 2, \dots).
\]

Hence \Olim $n^{-1}(1,0) \neq 0$. A necessary and sufficient condition for the validity of the equation 
\begin{equation}\label{eq:XII.2.24}
\Olim _ {\alpha_{n} \rightarrow \alpha} \alpha_{n} x = \alpha x
\end{equation}
is the so-called Archimedean axiom 
\begin{equation}\label{eq:XII.2.25}
\Olim _{n \to \infty} n ^{-1} x = 0 \quad \text{ for every } \quad x \geqq 0.
\end{equation}

This we see by the Jordan decomposition $y = y^{+} + y^{-}$. 

\begin{definition}\label{def:XII.2.Definition.3} $A$ subset $\{x_{\alpha}\}$ of a vector lattice $X$ is said to be bounded if there exist $y$ and $z$ such that $y \leqq x_{\alpha} \leqq z$ for all $x_{\alpha}$. $X$ is said to be complete if, for any bounded set $\{x_{\alpha}\}$ of $X$, $\sup\limits_{\alpha} x_{\alpha}$ and $\inf\limits_{\alpha} x_{\alpha}$ exist in $X$. Here $\sup\limits_{\alpha} x_{\alpha}$ is the least upper bound in the sense of the partial order in $X$, and $\inf\limits_{\alpha} x_{\alpha}$ is the largest lower bound in the sense of the partial order in $X$. A vector lattice $X$ is said to be $\sigma$-complete, if, for any bounded sequence $\{x_{n}\}$ of $X$, $\sup\limits_{n\geqq 1}x_n$ and $\inf\limits_{n\geqq 1}x_n$ exist in $X$. We define, in a $\sigma$-complete vector lattice $X$, 
\begin{equation}\label{eq:XII.2.26}
O _{n \to \infty} ^{- \lim} x _{n} = \inf\limits _{m} \left(\sup\limits _{n \geqq m} x _{n}\right), O _{n \to \infty} ^{- \lim} x _{n} = \sup\limits _{m} \left(\inf\limits _{n \geqq m} x _{n}\right).
\end{equation}
\end{definition}

\begin{proposition}\label{prop:XII.2.Proposition.7} \Olim $x_{n} = x$ iff \Olim $x_{n} = O$-lim $x_{n} = x$. 
\end{proposition}

\begin{proof}
 Suppose $|x - x_n| \leqq w_n, w_n \downarrow 0$. Then $x - w_n \leqq x_n \leqq x + w_n$ and so we obtain $O\text{- }\varlimsup\limits_{n \to \infty} (x - w_n) = x \leqq O\text{- }\varliminf\limits_{n \to \infty} x_n \leqq O\text{- }\varliminf\limits_{n \to \infty} (x + w_n) = x$, that is, $O\text{- }\varliminf\limits_{n \to \infty} x_n = x$. Similarly, we obtain $O\text{- }\varliminf\limits_{n \to \infty} x_n = x$. 

We next prove the sufficiency. Put $u_{n} = \sup\limits_{m\geqq n}x_{m}, v_{n} = \inf\limits_{m\geqq n}x_{m}, u_{n} - v_{n} = w_{n}$. Then, by hypothesis, $w_{n}\downarrow 0$. Also, by $x_{n}\leqq u_{n} = x + (u_{n} - x)\leqq x + (u_{n} - v_{n}) = x + w_{n}$ and $x_{n}\geqq x - w_{n}$ obtained similarly, we prove $|x - x_n|\leqq w_n$. Thus $\Olim_{n\to \infty}x_n = x$. 
\end{proof}

\begin{proposition}\label{prop:XII.2.Proposition.8} In a $\sigma$-complete vector lattice $X$, $\alpha x$ is continuous in $\alpha$, $x$ with respect to \Olim. 
\end{proposition}

\begin{proof}
 We have $|\alpha x - \alpha_n x_n| \leqq |\alpha x - \alpha x_n| + |\alpha x_n - \alpha_n x_n| = |\alpha| |x - x_n| + |\alpha - \alpha_n| |x_n|$. Hence, if $\Olim_{n \to \infty} x_n = x$, $\lim\limits_{n \to \infty} \alpha_n = \alpha$, then the $\Olim$ of the first term on the right is $0$. Therefore, by putting $\sup\limits_{n \geqq 1} |x_n| = y$, $\sup\limits_{m \geqq n} |\alpha - \alpha_m| = \beta_n$, we have to prove $\Olim_{n \to \infty} \beta_n y = 0$. But, by $y \geqq 0$ and $\beta_n \downarrow 0$, we see that $\Olim_{n \to \infty} \beta_n y = z$ exists and $\Olim_{n \to \infty} 2^{-1} \beta_n y = 2^{-1} z$. Since there exists, for any $n$, an $n_0$ such that $\beta_{n_0} \leqq 2^{-1} \beta_n$, we must have $z = 2^{-1} z$, that is, $z = 0$. 
\end{proof}

\begin{proposition}\label{prop:XII.2.Proposition.9} In a $\sigma$-complete vector lattice $X$, \Olim $x_{n}$ exists iff 
\begin{equation}\label{eq:XII.2.27}
\underset {n, m \rightarrow \infty} {O - \lim} | x _{n} - x _{m} | = 0.
\end{equation}
\end{proposition}

\begin{proof}
 The necessity is clear from $|x_{n} - x_{m}| \leqq |x_{n} - x| + |x - x_{m}|$. If we set $|x_{n} - x_{m}| = y_{nm}$, then $O\text{- }\varlimsup\limits_{n \to \infty} x_{n} \leqq x_{m} + O\text{- }\varliminf\limits_{n \to \infty} y_{nm}, O\text{- }\varliminf\limits_{n \to \infty} x_{n} \geqq x_{m} - O\text{- }\varliminf\limits_{n \to \infty} y_{nm}$. Hence 
\[
0 \leqq O - \varlimsup\limits_{n \rightarrow \infty} x _{n} - O - \varliminf\limits_{n \rightarrow \infty} x _{n} \leqq O - \lim\limits _{m \rightarrow \infty} \left(O - \varlimsup\limits_{n \rightarrow \infty} y _{nm} - O - \varliminf\limits_{n \rightarrow \infty} y _{nm}\right) = 0. \]

Thus we have proved the sufficiency. 
\end{proof}

\begin{proposition}\label{prop:XII.2.Proposition.10} $A$ vector lattice $X$ is $\sigma$-complete iff every monotone increasing, bounded sequence $\{x_{n}\} \subseteq X$ has $\sup x_{n}$ in $X$. 
\end{proposition}

\begin{proof}
 We have only to prove the sufficiency. Let $\{z_n\}$ be any bounded sequence in $X$, and set $x_{n} = \sup\limits_{m\leqq n}z_{m}$. Then, by hypothesis, $\sup\limits_{n\geqq 1}x_{n} = z$ exists in $X$, and $z = \sup\limits_{n\geqq 1}z_{n}$. Similarly, we see that $\inf\limits_{n\geqq 1}z_n = \inf\limits_{n\geqq 1}\left(\inf\limits_{m\leqq n}z_m\right)$ also exists in $\dot{X}$. 
\end{proof}


\yossection{3}{\texorpdfstring{$B$}{B}-lattices and \texorpdfstring{$F$}{F}-lattices}
\label{sec:XII.3}

\begin{definitionu} $A$ real $B$-space ( $F$-space) $X$ is said to be a $B$-lattice ( $F$-lattice) if it is a vector lattice such that 
\begin{equation}\label{eq:XII.3.1}
| x | \leqq | y | \quad \text{ implies } \quad \| x \| \leqq \| y \|.
\end{equation}
\end{definitionu}

\begin{examplen} $C(S), L^p(S)$ are $B$-lattices by the natural partial order $x \leqq y$ which means $x(s) \leqq y(s)$ on $S(x(s) \leqq y(s))$ a.e. on $S$ for the case $L^p(S)$). $M(S, \mathfrak{B}, m)$ with $m(S) < \infty$ is an $F$-lattice by the natural partial order as in the case of $L^p(S)$. $A(S, \mathfrak{B})$ is a $B$-lattice by
$\left\| x \right\| = | x | (S) =$ the total variation of $x$ over $S.$


We have, by \eqref{eq:XII.3.1} and $|x| = |(|x|)|$, 
\begin{equation}\label{eq:XII.3.2}
\left\| x \right\| = \left\| (| x |) \right\|.
\end{equation}

In $A(S, \mathfrak{B})$, we have, moreover, 
\begin{equation}\label{eq:XII.3.3}
x \geqq 0, y \geqq 0 \quad \text{ imply } \quad \| x + y \| = \| x \| + \| y \|.
\end{equation}

S. \textsc{Kakutani} called a $B$-lattice satisfying \eqref{eq:XII.3.3} an abstract $L^1$-space. From \eqref{eq:XII.3.3} we have 
\begin{equation}\label{eq:XII.3.4}
| x | < | y | \quad \text{ implies } \quad \| x \| < \| y \|.
\end{equation}

The norm in $A(S, \mathfrak{B})$ is continuous in \Olim, that is 
\begin{equation}\label{eq:XII.3.5}
\Olim _{n \to \infty} x _{n} = x \quad \text{ implies } \quad \lim\limits _{n \to \infty} \| x _{n} \| = \| x \|.
\end{equation}

For, in $A(S, \mathfrak{B})$, $\Olim_{n \to \infty} x_n = x$ is equivalent to the existence of $y_n \in A(S, \mathfrak{B})$ such that $|x - x_n| (S) \leqq y_n(S)$ with $y_n(S) \downarrow 0$. It is easy to see that $M(S, \mathfrak{B}, m)$ satisfies \eqref{eq:XII.3.4} and \eqref{eq:XII.3.5}. 

\end{examplen}
\begin{proposition}\label{prop:XII.3.Proposition.1} $A$ $\sigma$-complete $F$-lattice $X$ satisfying \eqref{eq:XII.3.4} and \eqref{eq:XII.3.5} is a complete lattice. In particular, $A(S, \mathfrak{B})$ and $L^p(S)$ are complete lattices. 
\end{proposition}

\begin{proof}
 Let $\{x_{\alpha}\} \subseteq X$ be bounded. We may assume $0 \leqq x_{\alpha} \leqq x$ for all $\alpha$, and we shall show that $\sup\limits_{\alpha} x_{\alpha}$ exists. Consider the totality $\{z_{\beta}\}$ of 

$z_{\beta}$ obtained as the sup of a finite number of $x_{\alpha}$'s: $z_{\beta} = \bigvee\limits_{j=1}^{n} x_{\alpha_j}$. Set $\gamma = \sup\limits_{\beta} \|z_{\beta}\|$. Then there exists a sequence $\{z_{\beta_j}\}$ such that $\lim\limits_{j \to \infty} \|z_{\beta_j}\| = \gamma$. If we put $z_n = \sup\limits_{j \leqq n} z_{\beta_j}$, then $O - \lim\limits_{n \to \infty} z_n = w$ exists, and, by \eqref{eq:XII.3.5} and the definition of $\gamma$, we have $\|w\| = \gamma$. We shall prove that $w = \sup\limits_{\alpha} x_{\alpha}$. Suppose that $x_{\alpha} \vee w > w$ for a certain $x_{\alpha}$. Then, by \eqref{eq:XII.3.4}, $\|x_{\alpha} \vee w\| > \|w\| = \gamma$. But, by $x_{\alpha} \vee w = O - \lim\limits_{n \to \infty} (x_{\alpha} \vee z_n)$, $x_{\alpha} \vee z_n \in \{z_{\beta}\}$ and \eqref{eq:XII.3.5}, we have $\|x_{\alpha} \vee w\| = \lim\limits_{n \to \infty} \|x_{\alpha} \vee z_n\| \leqq \gamma$, which is a contradiction. Hence we must have $w \geqq x_{\alpha}$ for all $x_{\alpha}$. Let $x_{\alpha} \leqq u$ for all $x_{\alpha}$. Suppose $w \wedge u < w$.


Then, by \eqref{eq:XII.3.4}, $\|w \wedge u\| < \gamma$, contrary to the fact that $w \wedge u \geqq z_{\beta}$ for all $z_{\beta}$. Hence we must have $w = \sup x_{\alpha}$. 
\end{proof}

\begin{remarku} In $C(S)$, $\Olim_{n \to \infty} x_n = x$ does not necessarily imply $\slim\limits_{n \to \infty} x_n = x$. In $M(S, \mathfrak{B}, m)$, $\slim\limits_{n \to \infty} x_n = x$ does not necessarily imply $\Olim_{n \to \infty} x_n = x$. To see this, let $x_1(s), x_2(s), \ldots$ be the defining functions of the intervals of [$0$, $1$]: 
\[
\left[ 0, \frac {1}{2} \right], \left[ \frac {1}{2}, \frac {2}{2} \right], \left[ 0, \frac {1}{4} \right], \left[ \frac {1}{4}, \frac {2}{4} \right], \left[ \frac {2}{4}, \frac {3}{4} \right], \left[ \frac {3}{4}, \frac {4}{4} \right], \left[ 0, \frac {1}{8} \right], \left[ \frac {1}{8}, \frac {2}{8} \right], \dots .
\]
\end{remarku}

Then the sequence $\{x_{n}(s)\} \subseteq M([0,1])$ converges to $0$ asymptotically but does not converge to $0$ a.e., that is, we have $s\text{- }\lim\limits_{n\to \infty}x_n = 0$ but $O\text{- }\lim\limits_{n\to \infty}x_n = 0$ does not hold. 

\begin{proposition}\label{prop:XII.3.Proposition.2} Let $X$ be an $F$-lattice. Let a sequences $\{x_{n}\} \subseteq X$ satisfy $\slim\limits_{n \to \infty} x_{n} = x$. Then $\{x_{n}\}$ *-converges to $x$ relative uniformly. This means that, from any subsequence $\{y_{n}\}$ of $\{x_{n}\}$, we can choose a subsequence $\{y_{n(k)}\}$ and a $z \in X$ such that 
\begin{equation}\label{eq:XII.3.6}
\left| y _{n(k)} - x \right| \leqq k ^{-1} z \quad (k = 1, 2, \dots).
\end{equation}
\end{proposition}

Conversely, if $\{x_{n}\}^{*}$-converges relative uniformly to $x$, then \slim $x_{n} = x$. 

\begin{proof}
 We may restrict ourselves to the case $x = 0$. From $\lim\limits_{n \to \infty} \| y_n \| = 0$, we see that there exists a sequence $\{n(k)\}$ of positive integers such that $\| k y_{n(k)} \| \leqq k^{-2}$. Then \eqref{eq:XII.3.6} holds by taking $z = \sum\limits_{k=1}^{\infty} |k y_{n(k)}|$. Conversely, let condition \eqref{eq:XII.3.6} with $x = 0$ be satisfied. Then we have $\| y_{n(k)} \| \leqq \| k^{-1} z \|$ from $|y_{n(k)}| \leqq k^{-1} z$. Hence $s \cdot \lim\limits_{k \to \infty} y_{n(k)} = 0$. Therefore, there cannot exist a subsequence $\{y_n\}$ of $\{x_n\}$ such that $\lim\limits_{n \to \infty} \| y_n \| > 0$. 
\end{proof}

\begin{remarku} The above proposition is an abstraction of the fact, that an asymptotically convergent sequence of $M(S, \mathfrak{B}, m)$ with $m(S) < \infty$ contains a subsequence which converges $m$-a.e. 
\end{remarku}


\yossection{4}{A Convergence Theorem of Banach}
\label{sec:XII.4}

This theorem is concerned with the almost everywhere convergence of a sequence of linear operators whose ranges are measurable functions. See S. \textsc{Banach} \cite{banach10}. A lattice-theoretic formulation of the theorem reads as follows (K. \textsc{Yosida} \cite{yosida358}): 

\begin{theoremu} Let $X$ be a real $B$-space with the norm $\|\|$ and $Y$ a $\sigma$-complete $F$-lattice with the quasi-norm $\|\|_1$ such that 
\begin{equation}\label{eq:XII.4.1}
\Olim _{n \to \infty} y _{n} = y \quad \text{ implies } \quad \lim\limits _{n \to \infty} \| y _{n} \| _{1} = \| y \| _{1}.
\end{equation}
\end{theoremu}

Let $\{T_n\}$ be a sequence of bounded linear operators $\in L(X, Y)$. Suppose that 
\begin{equation}\label{eq:XII.4.2}
\Olim _{n \to \infty} | T _{n} x | \text{ exists for those } x ^{\prime} s \in X \text{ which form a set }
\end{equation}

$G$ of the second category. 

Then, for any $x \in X$, $O - \varliminf\limits_{n \to \infty} T_n x$ and $O - \varliminf\limits_{n \to \infty} T_n x$ both exist and the (not necessarily linear) operator $\tilde{T}$ defined by 
\begin{equation}\label{eq:XII.4.3}
\tilde {T} x = \left(O - \varlimsup\limits_{n \rightarrow \infty} T _{n} x\right) - \left(O - \varliminf\limits_{n \rightarrow \infty} T _{n} x\right)
\end{equation}
is continuous as an operator defined on $X$ into $Y$. 

\begin{remarku} The space $M(S, \mathfrak{B}, m)$ with $m(S) < \infty$ satisfies \eqref{eq:XII.4.1} if we write $y_1 \leqq y_2$ when and only when $y_1(s) \leqq y_2(s)$ $m$-a.e., the quasi-norm $\|y\|_1$ being defined by $\|y\|_1 = \int_S |y(s)|(1 + |y(s)|)^{-1} m(\dd s)$. Likewise $L^p(S, \mathfrak{B}, m)$ with $m(S) \leqq \infty$ also satisfies \eqref{eq:XII.4.1} by the same semi-order. 
\end{remarku}

\begin{proof}[Proof of the Theorem.]
Set $T_{n}x = y_{n}, y_{n}' = \sup\limits_{n\geqq m}|y_{m}|, y' = \sup\limits_{n\geqq 1}|y_{n}|$, and consider the operators $V_{n}x = y_{n}'$ and $Vx = y'$ defined at least on $G$ into $Y$. By ($23$) of the preceding section $2$, each $V_{n}$ is strongly continuous with $T_{k}$. Since $\lim\limits_{n\to \infty}\|V_n x - V x\|_1 = 0$ by \eqref{eq:XII.4.1}, we have $\lim\limits_{n\to \infty}\|k^{-1}V_n x\| = \|k^{-1}Vx\|$ and further $\lim\limits_{k\to \infty}\|k^{-1}Vx\|_1 = 0$. These are implied by the continuity in $\alpha, y$ of $\alpha y$ in the $F$-space $Y$. Hence for any $\varepsilon > 0$, 
\begin{equation}\label{eq:XII.4.4}
G \subseteq \bigcup\limits_{k=1}^{\infty} G_k \text{ where } G_k = \{x \in X; \sup\limits_{n\geqq 1}\|k^{-1}V_n x\|_1 \leqq \varepsilon\}.
\end{equation}

By the strong continuity of $V_{n}$, each $G_{k}$ is a strongly closed set of $X$. Thus some $G_{k_0}$ contains a sphere of $X$, in virtue of the hypothesis that $G$ is of the second category. That is, there exist an $x_0 \in X$ and a $\delta > 0$ such that $\|x_0 - x\| \leqq \delta$ implies $\sup\limits_{n \geqq 1} \|k_0^{-1} V_n x\|_1 \leqq \varepsilon$. Hence, by putting $z = x_0 - x$, we see that 
\[
\sup\limits _{n \geqq 1} \| k _{0} ^{-1} V _{n} z \| _{1} \leqq \sup\limits _{n \geqq 1} \| k _{0} ^{-1} V _{n} x _{0} \| _{1} + \sup\limits _{n \geqq 1} \| k _{0} ^{-1} V _{n} x \| _{1} \leqq 2 \varepsilon ,
\]
that is, since $V_{n}(k_{0}^{-1}z) = k_{0}^{-1}V_{n}z$, we have 
\[
\sup\limits _{n \geqq 1} \| V _{n} z \| _{1} \leqq 2 \varepsilon \text{ whenever } \| z \| \leqq \delta / k _{0}.
\]

This proves that $\slim\limits_{\|z\| \to 0} V_n \cdot z = 0$ uniformly in $n$. 

$G$ being dense in $X$, $V \cdot x$ is defined for all $x \in X$ and $V \cdot x$ is strongly continuous at $x = 0$ with $V \cdot 0 = 0$. Hence, by 
\[
| \tilde {T} \cdot x | \leqq 2 V \cdot x \quad \text{ and } \quad \| \tilde {T} x _{1} - \tilde {T} x _{2} \| _{1} \leqq \| \tilde {T} (x _{1} - x _{2}) \| _{1},
\]
we see that $\tilde{T} \cdot x$ is strongly continuous at every $x \in X$. 
\end{proof}

\begin{corollaryu} Under condition \eqref{eq:XII.4.1}, the set $G = \{x \in X; O - \lim\limits_{n \to \infty} T_n x \text{ exists }\}$ either coincides with $X$ or is a set of the first category. 
\end{corollaryu}

\begin{proof}
 Suppose the set $G$ be of the second category. Then, by our Theorem, the operator $\tilde{T}$ is a strongly continuous operator defined on $X$ into $Y$. Thus $G = \{x \in X; \tilde{T}y = 0\}$ is strongly closed in $X$. Moreover, $G$ is a linear subspace of $X$. Hence $G$ must coincide with $X$. Otherwise, $G$ would be non dense in $X$. 
\end{proof}


\yossection{5}{The Representation of a Vector Lattice as Point Functions}
\label{sec:XII.5}

Let a vector lattice $X$ contain a ``unit'' $I$ with the properties: 

$I > 0$, and, for any $f \in X$, there exists an $\alpha > 0$ such 
\begin{equation}\label{eq:XII.5.1}
\text{ that } - \alpha I \leqq f \leqq \alpha I.
\end{equation}

For such a vector lattice $X$, we can give an analogous representation similar to the representation of a normed ring as point functions. 

An element $f \in X$ is called ``nilpotent'' if $n |f| \leqq I$ ( $n = 1, 2, \ldots$). The set $R$ of all nilpotent elements $f \in X$ is called the ``radical'' of $X$. By ($20$) of Chapter~\ref{chap:XII}, \ref{sec:XII.2}, $R$ constitutes a linear subspace of $X$. Moreover, $R$ is an ``ideal'' of $X$ in the sense that 
\begin{equation}\label{eq:XII.5.2}
f \in R \text{ and } | g | \leqq | f | \text{ imply } g \in R.
\end{equation}

\begin{lemmau} Let $X_{1}$ and $X_{2}$ be vector lattices. A linear operator $T$ defined on $X_{1}$ onto $X_{2}$ is called a lattice homomorphism if 
\begin{equation}\label{eq:XII.5.3}
T (x \underset {\wedge} {\vee} y) = (T x) \underset {\wedge} {\vee} (T y).
\end{equation}
\end{lemmau}

Then \eqref{eq:XII.5.3} implies that $N = \{x \in X_1; Tx = 0\}$ is an ideal of $X_1$. Conversely, if a linear subspace $N$ of $X_1$ is an ideal of $X_1$, then the mapping $x \to Tx = \overline{x} = (\text{ the residue class mod } N \text{ containing } x)$ satisfies \eqref{eq:XII.5.3} by defining $\overline{x}_{\wedge}^{\vee}\overline{y}$ through $\overline{x}_{\wedge}^{\vee}\overline{y}$. 

\begin{proof}
 Let $T$ be a lattice homomorphism. Let $x \in N$ and $|y| \leqq |x|$. Then, by $T(|x|) = T(x \vee -x) = (Tx) \vee (T(-x)) = 0$, we obtain, $0 \leqq Ty^{+} = T(y^{+} \wedge |x|) = Ty^{+} \wedge T|x| = 0$ and so $y^{+} \in N$. Similarly we obtain $y^{-} \in N$ and thus $y = y^{+} + y^{-} \in N$. 


Next let a linear subspace $N = \{x \in X_1; Tx = 0\}$ be an ideal of $X_1$. If $\overline{y} = \overline{z}$, i.e., if $y - z \in N$, then, by \eqref{eq:XII.2.23} in Chapter~\ref{chap:XII}, \ref{sec:XII.2}, we have 
\[
\left| x _{\wedge} ^{\vee} y - x _{\wedge} ^{\vee} z \right| \leqq | y - z | \in N
\]
so that the residue class $(x \stackrel{\vee}{\wedge} y)$ is determined independently of the choice of the representative elements $x$, $y$ from the residue classes $\overline{x}, \overline{y}$, respectively. Hence we can put $\overline{x} \stackrel{\vee}{\wedge} \overline{y} = \overline{x} \stackrel{\vee}{\wedge} y$ by defining $\overline{z} \geqq \overline{0}$ through $z \vee 0 = \overline{z}$.
\end{proof}

\begin{remarku} The above Lemma may be phrased as follows. Let $N$ be a linear subspace of a vector lattice. Then the linear-congruence $a \equiv b \pmod{N}$ is also a lattice-congruence: 
\[
a \equiv b, a ^{\prime} \equiv b ^{\prime} (\mathrm{mod} N) \quad \text{ implies } \quad a \underset {\wedge} {\vee} b \equiv a ^{\prime} \underset {\wedge} {\vee} b ^{\prime} (\mathrm{mod} N),
\]
\end{remarku}

iff $N$ is an ideal of $X$. 

Now an ideal $N$ is called ``non-trivial'' if $N \neq \{0\}$, $X$. A non-trivial ideal $N$ is called ``maximal'' if it is contained in no other ideal $\neq X$. Denote by $M$ the set of all maximal ideals $N$ of $X$. The residual class $X$/$N$ of $X$ mod any ideal $N \in M$ is ``simple'', that is, $X$/$N$ does not contain non-trivial ideals. It will be proved below that a simple vector lattice with a unit is linear-lattice isomorphic to the vector lattice of real numbers, the non-negative elements and the unit $I$ being represented by non-negative numbers and the number $1$. We denote by $f(N)$ the real number which corresponds to $f \in X$ by the linear-lattice-homomorphism $X \to X/N$, $N \in M$. 

After these preliminaries we may state 

\begin{theorem}\label{thm:XII.5.Theorem.1} The radical $R$ coincides with the intersection ideal $\bigcap\limits_{N \in \mathfrak{M}} N$. 
\end{theorem}

\begin{proof}
 The first step. Let $X$ be a simple vector lattice with a unit $I$. Then we must have $X = \{\alpha I; -\infty < \alpha < \infty\}$. 


\begin{proof}
 $X$ does not contain a nilpotent element $f \neq 0$, for otherwise $X$ would contain a non-trivial ideal $N = \{g; |g| \leqq \eta |f|$ with some $\eta < \infty\}$. Hence, by \eqref{eq:XII.5.1}, $X$ satisfies the Archimedean axiom: 
\begin{equation}\label{eq:XII.5.4}
\underset {n \rightarrow \infty} {\text{ order-lim }} n ^{-1} | x | = 0 \quad \text{ for all } \quad x \in X. \end{equation}

Suppose there exists an $f_0 \in X$ such that $f_0 \neq \gamma I$ for any real number $\gamma$. Let 
\[
\alpha = \inf\limits _ {f _{0} \leqq \alpha^{\prime} I} \alpha^{\prime}, \beta = \sup\limits _ {\beta^{\prime} I \leqq f _{0}} \beta^{\prime}.
\]

Then, by \eqref{eq:XII.5.4}, $\beta I \leqq f_0 \leqq \alpha I$ and $\beta < \alpha$. Hence $(f_0 - \delta I)^+ \neq 0$, $(f_0 - \delta I)^- \neq 0$ for $\beta < \delta < \alpha$. Thus, by $x^+ \wedge (-x^-) = 0$, the set $N_0 = \{g; |g| \leqq \eta(f_0 - \delta I)^+$ with some $\eta < \infty\}$ is a non-trivial ideal, contrary to the hypothesis. 
\end{proof}

The second step. For any non-trivial ideal $N_0$, there exists a maximal ideal $N_1$ containing $N_0$. 

\begin{proof}
 Let $\{N_0\}$ be the totality of non-trivial ideals containing $N_0$. We order the ideals of $\{N_0\}$ by inclusion relation, that is, we denote $N_{\alpha_1} \leqq N_{\alpha_2}$ if $N_{\alpha_1}$ is a subset of $N_{\alpha_2}$. Suppose that $\{N_\alpha\}$ is a linearly ordered subset of $\{N_0\}$ and set $N_\beta = \bigcup\limits_{N_\alpha \in \{N_\alpha\}} N_\alpha$. We shall then show that $N_\beta$ is an upper bound of $\{N_\alpha\}$. For, if $x, y \in N_\beta$, there exist ideals $N_{\alpha_1}$ and $N_{\alpha_2}$ such that $x \in N_{\alpha_1}$ and $y \in N_{\alpha_2}$. Since $\{N_\alpha\}$ is linearly ordered, $N_{\alpha_1} \subseteq N_{\alpha_2}$ (or $N_{\alpha_1} \supseteq N_{\alpha_2}$) and so $x$ and $y$ both belong to $N_{\alpha_2}$. This proves that $(\gamma x + \delta y) \in N_{\alpha_1} \subseteq N_\beta$ and that $|z| \leqq |x|$ implies $z \in N_{\alpha_1} \subseteq N_\beta$. Since the unit $I$ is contained in no $N_\alpha$, $I$ is not contained in $N_\beta$. Therefore $N_\beta$ is a non-trivial ideal containing every $N_\alpha$, i.e., $N_\beta$ is an upper bound of $\{N_\alpha\}$. Thus, by Zorn's Lemma, there exists at least one maximal ideal containing $N_0$. 
\end{proof}

The third step. $R \subseteq \bigcap\limits_{N \in \mathfrak{M}} N$. Let $f > 0$ and $nf \leqq I (n = 1, 2, \ldots)$. Then; for any $N \in \mathfrak{M}$, we have $nf(N) \leqq I(N) = 1 (n = 1, 2, \ldots)$, and hence $f(N) = 0$, that is, $f \in N$. 

The fourth step. $R \supseteq \bigcap\limits_{N \in \mathfrak{M}} N$. Let $f > 0$ be not nilpotent. Then we have to show that there exists an ideal $N \in \mathfrak{M}$ such that $f \bar{\in} N$. This may be proved as follows. 

Since $f > 0$ is not nilpotent, there exists an integer $n$ such that $nf \not\leqq I$. We may assume that $nf \not\leqq I$, since otherwise $f \not\in N$ for any $N \in \mathfrak{M}$ and so we have nothing to prove. Thus suppose $p = I - (n \cdot f) \cap I > 0$. Then, for any positive integer $m$, we do not have $m \cdot p \geqq I$. If otherwise, we would have $m^{-1}I \leqq I - (n \cdot f) \cap I$ and hence 
\[
(n \cdot f) \wedge I = (n \cdot f) \wedge (1 - m ^{-1}) I.
\]

Thus, by \eqref{eq:XII.2.6} in Chapter~\ref{chap:XII}, \ref{sec:XII.2}, 
\[
(n \cdot f - (1 - m ^{-1}) I) \wedge m ^{-1} I = (n \cdot f - (1 - m ^{-1}) I) \wedge 0 \leqq 0,
\]
and so, by the distributivity of the vector lattice, 
\[
0 = \{(n \cdot f - (1 - m ^{-1}) I) \wedge m ^{-1} I \} \vee 0 = (n \cdot f - (1 - m ^{-1}) I) ^{+} \wedge m ^{-1} I,
\]
that is, $(n \cdot f - (1 - m^{-1}) I)^{+} \wedge I = 0$. Put $b = (n \cdot f - (1 - m^{-1}) I)^{+}$ and assume that $b > 0$. By hypothesis \eqref{eq:XII.5.1}, we have $b < \alpha I$ with some $\alpha > 1$. Then $0 < b = b \wedge \alpha I$ and so $0 < (\alpha^{-1} b) \wedge I \leqq b \wedge I$, contrary to $b \wedge I = 0$. Therefore $b = 0$, i.e., $n \cdot f \leqq (1 - m^{-1}) I$. This contradicts the fact that $n \cdot f \not\leqq I$. Hence the set $N_{0} = \{g; |g| \leqq \eta |p|$ for some $\eta < \infty\}$ is a non-trivial ideal. $N_{0}$ is contained in at least one maximal ideal $N$, by the second step. Then $0 = p(N) = 1 - (n \cdot f(N)) \wedge 1$ which shows that $f(N) > 0$, that is, $f \not\in N$. 

We have thus proved our Theorem~\ref{thm:XII.5.Theorem.1}.
\end{proof} 

The vector lattice $\overline{X} = X / R$ is again a vector lattice with a unit $\overline{I}$. By Theorem~\ref{thm:XII.5.Theorem.1}, the intersection ideal $\bigcap\limits_{\overline{N}} \overline{N}$ of all the maximal ideals $\overline{N}$ of $\overline{X}$ is the zero ideal and $\overline{X}$ contains no nilpotent element $\neq 0$. Hence $\overline{X}$ satisfies the Archimedean axiom 
\begin{equation}\label{eq:XII.5.5}
\underset {n \uparrow \infty} {\text{ order-limit }} n ^{-1} | \bar {f} | = 0 \quad \text{ for all } \quad \bar {f} \in \overline {{X}}.
\end{equation}

Let $\bar{N}$ be any maximal ideal of $\bar{X}$. Then the factor space $\bar{X}/\bar{N}$ is a simple vector lattice, and so, by the first step in the proof of Theorem~\ref{thm:XII.5.Theorem.1}, $\bar{X}/\bar{N}$ is linear-lattice-isomorphic to the vector lattice of real numbers; the non-negative elements and the unit being represented by non-negative numbers and $1$. We denote by $\overline{f}(\overline{N})$ the real number which corresponds to $\overline{f}$ by this homomorphism $\overline{X} \to \overline{X}/\overline{N}$. We also denote by $\overline{\mathfrak{M}}$ the set of all maximal ideals of $\overline{X}$. We have thus 

\begin{theorem}\label{thm:XII.5.Theorem.2} By the correspondence $\overline{f} \to \overline{f}(\overline{N})$, $\overline{X}$ is linear-lattice-isomorphically mapped on the vector lattice $F(\mathfrak{M})$ of real-valued bounded functions on $\overline{\mathfrak{M}}$ such that \begin{enumerate}[(i)]
\item $|\overline{f}| \to |\overline{f}(\overline{N})|$,
\item $\overline{I}(\overline{N}) \equiv 1$ on $\overline{\mathfrak{M}}$
\item $F(\overline{\mathfrak{M}})$ separates the points of $\overline{\mathfrak{M}}$ in the sense that 

\begin{equation}\label{eq:XII.5.6}
\begin{minipage}[b]{0.92\linewidth}
for two different points $\overline{N}_1, \overline{N}_2$ of $\overline{\mathfrak{M}}$, there exists at least one $\overline{f} \in \overline{X}$ such that $\overline{f}(\overline{N}_1) \neq \overline{f}(\overline{N}_2)$.
\end{minipage}
\end{equation}
\end{enumerate}
\end{theorem}


\begin{remarku} We introduce a topology in $\overline{M}$ by calling the sets of the form 
\begin{align}
\{\bar {N}& \in \overline {{\mathfrak {M}}}; | \bar {f} _{i} (\bar {N}) - \bar {f} _{i} (\bar {N} _{0}) | < \varepsilon_{i} (i = 1, 2, \dots , n), \notag \\
\text{ where } - \bar {I}& \leqq \bar {f} \leqq \bar {I} \quad \text{ for all } i \} \notag
\end{align}
\end{remarku}

neighbourhoods of $N_0$. Then $\overline{\mathfrak{M}}$ is compact since it may be identified with a closed subset of a topological product (of the same potency as the cardinal number of the set of elements $\bar{f} \in X$ which satisfy $-\bar{I} \leqq \bar{f} \leqq \bar{I}$) of the closed intervals [---$1$, $1$]. The proof is entirely similar to the case of the set of all maximal ideals of a normed ring in Chapter~\ref{chap:XI}, \ref{sec:XI.2}. Moreover, each function $\bar{f}(\overline{N}) \in F(\overline{\mathfrak{M}})$ is continuous on the compact space $\overline{\mathfrak{M}}$ topologized in this way. Thus, by the Kakutani-Krein theorem in Chapter~\ref{chap:0}, \ref{sec:0.2}, we see that $F(\overline{\mathfrak{M}})$ is dense in the $B$-space $C(\overline{\mathfrak{M}})$. The above two theorems are adapted from K. \textsc{Yosida}-M. \textsc{Fukamiya} \cite{yosida359}. Cf. also S. \textsc{Kakutani} \cite{kakutani128} and M. \textsc{Krein}-S. \textsc{Krein} \cite{krein166}. 


\yossection{6}{The Representation of a Vector Lattice as Set Functions}
\label{sec:XII.6}

Let $X$ be a $\sigma$-complete vector lattice. Choose any positive element $x$ of $X$ and call it a ``unit'' of $X$ and write $1$ for $x$; when not ambiguous we also write $\alpha$ for $\alpha \cdot 1$. A non-negative element $e \in X$ is called a ``quasi-unit'' if $e \wedge (1 - e) = 0$. A finite linear combination $\sum\limits_{i} \alpha_{i} e_{i}$ of quasi-units $e_{i}$ is called a ``step-element'', and we call the element $y \in X$ ``absolutely continuous'' (with respect to the unit $1$) if $y$ can be expressed as the \Olim of a sequence of step-elements. An element $z \in X$ is called ``singular'' (with respect to the unit $1$) if $|z| \wedge 1 = 0$. 

We shall give an abstract formulation of the Radon-Nikodym theorem in integration theory. 

\begin{theoremu} Any element of $X$ is uniquely expressed as the sum of an absolutely continuous element and a singular element. 
\end{theoremu}

\begin{proof}
 The first step. If $f > 0$ and $f \wedge 1 \neq 0$, then there exist a positive number $\alpha$ and a quasi-unit $e_{\alpha} \neq 0$ such that $f \geqq \alpha e_{\alpha}$. We can, 

in fact, take 
\begin{equation}\label{eq:XII.6.1}
e _{\alpha} = \bigvee\limits_{n \geqq 1} \left\{n \left(\alpha^{-1} f - \alpha^{-1} f \wedge 1\right) \wedge 1 \right\}. \end{equation}


\begin{proof}
 Put $y_{\alpha} = \alpha^{-1}f - \alpha^{-1}; \wedge 1$. Then we obtain 
\[
2 e _{a} \wedge 1 = \left\{\bigvee\limits_{n \geqq 1} (2 n y _{a} \wedge 2) \right\} \wedge 1 = e _{\alpha},
\]
and so $e_{\alpha}$ is a quasi-unit. We obtain $f \geqq \alpha e_{\alpha}$ from 
\[ \begin{aligned} n y _{\alpha} \wedge 1 & = n \alpha^{-1} f \wedge [ 1 + n (\alpha^{-1} f \wedge 1) ] - n (\alpha^{-1} f \wedge 1) \\ & \leqq (n + 1) \alpha^{-1} f \wedge (n + 1) - n (\alpha^{-1} f \wedge 1) \leqq \alpha^{-1} f \wedge 1 \leqq \alpha^{-1} f. \end{aligned}
\]

If we can show that $y_{\alpha} \wedge 1 > 0$ for some $\alpha > 0$, then $e_{\alpha} > 0$ for such $\alpha$. Suppose that such positive $\alpha$ does not exist. Then, for every $\alpha$ with $0 < \alpha < 1$, we have 
\[
\alpha^{-1} (\alpha^{-1} f - \alpha^{-1} f \wedge 1) \wedge \alpha^{-1} = 0.
\]

Hence $(f - f \wedge \alpha) \wedge 1 = 0$ and, letting $\alpha \downarrow 0$, we obtain $f \wedge 1 = 0$, contrary to the hypothesis $f \wedge 1 \neq 0$. 
\end{proof}

The second step. Let $f \geqq 0$, and $f \geqq \alpha e$ where $\alpha > 0$ and $e$ is a quasi-unit. Then, for $0 < \alpha' < \alpha$, $e_{\alpha'} \geqq e$ and $f \geqq \alpha' e_{\alpha'}$ where $e_{\alpha'}$ is defined by \eqref{eq:XII.6.1}. 

\begin{proof}
 For the sake of simplicity, we assume that $\alpha = 1$. For $0 < \delta < 1$ 
\[ \begin{aligned}
&\frac {f}{1 - \delta} + \frac {e}{1 - \delta} \wedge 1  = \left(\frac {f}{1 - \delta} + \frac {e}{1 - \delta}\right) \wedge \left(1 + \frac {f}{1 - \delta}\right) \geqq \left(\frac {f}{1 - \delta} + \frac {e}{1 - \delta}\right) \wedge \left(1 + \frac {e}{1 - \delta}\right) \\
&= \frac {f}{1 - \delta} \wedge 1 + \frac {e}{1 - \delta}.
\end{aligned}
\]

Since $e$ is a quasi-unit, we have $2e \wedge 1 = e$, $e \wedge 1 = e$. Hence $me \wedge 1 = e$ when $m \geqq 1$. Thus, by $1 < (1 - \delta)^{-1}$, we have $(1 - \delta)^{-1}e \wedge 1 = e$. Therefore, we obtain, from the above inequality, 
\[
\frac {\delta}{1 - \delta} e = \frac {e}{1 - \delta} - \frac {e}{1 - \delta} \wedge 1 \leqq \frac {f}{1 - \delta} - \frac {f}{1 - \delta} \wedge 1 = y _{1- \delta},
\]
and so, by \eqref{eq:XII.6.1}, $e \leqq e_{1- \delta}$. 
\end{proof}

The third step. The set of all quasi-units constitutes a Boolean algebra, that is, if $e_{1}$ and $e_{2}$ are quasi-units, then $e_{1} \vee e_{2}$ and $e_{1} \wedge e_{2}$ are also quasi-units and $0 \leqq e_{i} \leqq 1$. The quasi-unit $(1 - e)$ is the complement of $e$, and $0$,$1$ are the least and the greatest elements, respectively in the totality of the quasi-units. 

\begin{proof}
 The condition $e \wedge (1 - e) = 0$ is equivalent to $2e \wedge 1 = e$. Hence, if $e_1, e_2$ are quasi-units, then 
\begin{align}
2 \left(e _{1} \wedge e _{2}\right) \wedge 1& = \left(2 e _{1} \wedge 1\right) \wedge \left(2 e _{2} \wedge 1\right) = e _{1} \wedge e _{2}, \notag \\
2 \left(e _{1} \vee e _{2}\right) \wedge 1& = \left(2 e _{1} \wedge 1\right) \vee \left(2 e _{2} \wedge 1\right) = e _{1} \vee e _{2}, \notag
\end{align}
so that $e_1 \wedge e_2$ and $e_1 \vee e_2$ are also quasi-units. 
\end{proof}

The fourth step. Let $f > 0$, and set $\overline{f} = \sup \beta e_{\beta}$ where sup is taken for all positive rational numbers $\beta$. Because of the third step, the sup of a finite number of elements of the form $\beta_i e_{\beta_i}$ is a step-element. Hence $\overline{f}$ is absolutely continuous with respect to the unit element $1$. We have to show that $g = f - \overline{f}$ is singular with respect to the unit element $1$. Suppose that $g$ is not singular. Then, by the first step, there exist a positive number $\alpha$ and a quasi-unit $e$ such that $g \geqq \alpha e$. Hence $f \geqq \alpha e$, and so, by the second step, there exists, for $0 < \alpha_1 < \alpha$, a quasi-unit $e_{a_1} \geqq e$ such that $f \geqq \alpha_1 e_{\alpha_1}$. We may assume that $\alpha_1$ is a rational number. Thus $\overline{f} \geqq \alpha_1 e_{\alpha_1}$ and so $f = \overline{f} + g \geqq 2\alpha_1 e_{\alpha_1}$. Again, by the second step, there exists, for $0 < \alpha_1' < \alpha_1$, a quasi-unit $e_{2\alpha_1'} \geqq e_{\alpha_1}$ such that $f \geqq 2\alpha_1' e_{2\alpha_1'}$. We may assume that $2\alpha_1'$ is a rational number so that $\overline{f} \geqq 2\alpha_1' e_{2\alpha_1'}$. Hence $f = \overline{f} + g \geqq 3\alpha_1' e$. Repeating the process, we can prove that, for any rational number $\alpha_n$ with $0 < \alpha_n < \alpha$, 
\[
f \geqq (n + 1) \alpha_{n} e \quad (n = 1, 2, \dots).
\]

If we take $\alpha_{n} \geqq \alpha / 2$, we have $(n + 1)\alpha_{n}e \geqq 2^{-1}n\alpha e$. Hence $f \geqq n\alpha e$ ( $n = 1, 2, \ldots$) with $\alpha > 0, e > 0$. This is a contradiction. For, by the $\sigma$-completeness of $X$, the Archimedean axiom holds in $X$. 

The fifth step. Let $f = f^{+} + f^{-}$ be the Jordan decomposition of a general element $f \in X$. Applying the fourth step to $f^{+}$ and $f^{-}$ separately, we see that $f$ is decomposed as the sum of an absolutely continuous element and a singular element. The uniqueness of the decomposition is proved if we can show that an element $h \in X$ is $= 0$ if $h$ is absolutely continuous as well as singular. But, since $h$ is absolutely continuous, we have $h = \Olim_{n \to \infty} h_{n}$ where $h_{n}$ are step-elements. $h_{n}$ being a step-element, there exists a positive number $\alpha_{n}$ such that $|h_{n}| \leqq \alpha_{n} \cdot 1$. Since $h$ is singular, we have $|h| \wedge |h_{n}| = 0$. Therefore $|h| = |h| \wedge |h| = \Olim_{n \to \infty} (|h| \wedge |h_{n}|) = 0$. 
\end{proof}

Application to the Radon-Nikodym theorem. Consider the case $X = A(S, \mathfrak{B})$. We know already (Proposition~\ref{prop:XII.3.Proposition.1} in Chapter~\ref{chap:XII}, \ref{sec:XII.3}) that it is a complete lattice, and that in $A(S, \mathfrak{B})$, 
\begin{equation}\label{eq:XII.6.2}
x^{+}(B) = \sup\limits_{B^{\prime}\subseteq B}x(B^{\prime}) = \text{ the positive variation } \overline{V} (x;B) \text{ of } x \text{ on } B.
\end{equation}

We shall prepare a 


\begin{propositionu} In $A(S, \mathfrak{B})$, let $x > 0$, $z \geqq 0$ in such a way that $x \wedge z = 0$. Then there exists a set $B_0 \in \mathfrak{B}$ such that $x(B_0) = 0$ and $z(S - B_0) = 0$. 
\end{propositionu}


\begin{proof}
 Since $x \wedge z = (x - z)^{-} + z$, we know, from \eqref{eq:XII.6.2}, that 
\begin{equation}\label{eq:XII.6.3}
(x \wedge z) (B) = \inf\limits _ {B ^{\prime} \subseteq B} (x - z) \left(B ^{\prime}\right) + z (B) = \inf\limits _ {B ^{\prime} \subseteq B} \left[ x \left(B ^{\prime}\right) + z \left(B - B ^{\prime}\right) \right]. \end{equation}

Hence, by the hypothesis $x \wedge z = 0$, we have 
\[
\inf\limits _ {B \in \mathfrak {B}} [ x (B) + z (S - B) ] = (x \wedge z) (S) = 0.
\]

Hence, for any $\varepsilon > 0$, there exists a $B_{\varepsilon} \in \mathfrak{B}$ such that $x(B_{\varepsilon}) \leqq \varepsilon$, $z(S - B_{\varepsilon}) \leqq \varepsilon$. Put $B_0 = \bigwedge\limits_{k \geqq 1} \left( \bigvee\limits_{n \geqq k} B_2^{-n} \right)$. Then, by the $\sigma$-additivity of $x(B)$ and $z(B)$, 
\begin{align}
0& \leqq x (B _{0}) = \lim\limits _{k \rightarrow \infty} x \left(\bigvee\limits_{n \geqq k} B _ {2 ^{-n}}\right) \leqq \lim\limits _{k \rightarrow \infty} \sum\limits_{n=k} ^{\infty} 2 ^{-n} = 0, \notag \\
0& \leqq z (S - B _{0}) = \lim\limits _{k \rightarrow \infty} z \left(S - \bigvee\limits_{n \geqq k} B _ {2 ^{-n}}\right) \leqq \lim\limits _{k \rightarrow \infty} z (S - B _ {2 ^{-k}}) = 0. \notag
\end{align}
\end{proof}

\begin{corollaryu} Let $e$ be a quasi-unit with respect to $x > 0$ in $A(S, \mathfrak{B})$. Then there exists a set $B_{1} \in \mathfrak{B}$ such that 
\begin{equation}\label{eq:XII.6.4}
e (B) = x \left(B \cap B _{1}\right) \quad \text{ for all } \quad B \in \mathfrak {B}.
\end{equation}
\end{corollaryu}

\begin{proof}
 Since $(x - e) \wedge e = 0$, there exists a set $B_0 \in \mathfrak{B}$ such that $e(B_0) = 0$, $(x - e)(S - B_0) = 0$. Hence $e(S - B_0) = x(S - B_0) = e(S)$, and so $e(B) = x(B - B_0) = x(B \cap B_1)$ where $B_1 = S - B_0$. 
\end{proof}

We are now able to prove the Radon-Nikodym theorem in integration theory. In virtue of the above Corollary, a quasi-unit $e$ with respect to $x > 0$ is, in $A(S, \mathfrak{B})$, a contracted measure $e(B) = x(B \cap B_e)$. Hence a step-element in $A(S, \mathfrak{B})$ is the integral of the form: 
\[
\sum\limits_{i} \lambda_{i} \int_{B \cap B i} x (\dd s),
\]
that is, indefinite integral of a step function (= finitely-valued function). Hence an absolutely continuous element in $A(S, \mathfrak{B})$ is an indefinite integral with respect to the measure $x(B)$. The Proposition above says that a singular element $g$ (with respect to the unit $x$) is associated with a set $B_0 \in \mathfrak{B}$ such that $x(B_0) = 0$ and $g(B) = g(B \cap B_0)$ for all $B \in \mathfrak{B}$. Such a measure $g(B)$ is the so-called singular measure (with respect to $x(B)$). Thus any element $f \in A(S, \mathfrak{B})$ is expressed as the sum of an indefinite integral (with respect to $x(B)$) and a measure $g(B)$ which is singular (with respect to $x(B)$). This decomposition is unique. The obtained result is exactly the Radon-Nikodym theorem. 

\begin{remarku} The above Theorem is adapted from K. \textsc{Yosida} \cite{yosida345}. Cf. F. \textsc{Riesz} \cite{rickart277}, H. \textsc{Freudenthal} \cite{freudenthal67} and S. \textsc{Kakutani} \cite{kakutani129}. For further references, see G. \textsc{Birkhoff} \cite{birkhoff17}. 
\end{remarku}
